% swift-closures-deep.tex — closure syntax, escaping vs non-escaping,
% withoutActuallyEscaping, @autoclosure, capture semantics (variables vs capture
% lists), the loop trap, closures as reference types, @Sendable, trailing closures,
% function types. Retain cycles are on ios-swift/arc-ownership.tex — only referenced.
% Sources (checked 2026-10-06):
%   docs.swift.org/swift-book Closures (capture semantics, escaping, autoclosures; optional
%     closure parameters are implicitly escaping) · Expressions › Capture Lists
%   developer.apple.com/documentation/swift/withoutactuallyescaping(_:do:) (runtime check that
%     the closure did not escape)
%   swift-evolution SE-0269 (implicit self, 5.3) · SE-0365 (implicit self after
%     guard let self, 5.8) · SE-0279 (multiple trailing closures, 5.3) · SE-0302 (Sendable /
%     @Sendable closures, 5.5) · SE-0430 (sending, 6.0) — checked against the evolution index
% @labels: area=swift kind=concept platform=apple topic=language,memory discussion=153
% @tags: closures, capture-list, escaping, non-escaping, autoclosure, withoutactuallyescaping, weak-self, sendable-closure, trailing-closures, retain-cycle, loop-capture
\documentclass[8pt]{extarticle}
\usepackage{printup-sheet}
\usepackage{array}

\lstdefinelanguage{SwiftSheet}{
  morekeywords={protocol,class,final,struct,enum,func,var,let,init,case,for,while,
    if,else,return,guard,self,nil,private,true,false,in,weak,unowned,async,await,
    throws,rethrows,try,static,some,Self,String,Int,Bool,Void,Data},
  sensitive=true, morecomment=[l]{//}, morecomment=[s]{/*}{*/}, morestring=[b]"}

% One \hbox per character defeats the mono font's ligatures (?? && ===).
\makeatletter
\DeclareRobustCommand\op[1]{\texttt{\@tfor\@c:=#1\do{\hbox{\@c}}}}
\makeatother

\tikzset{
  lbl/.style={font=\tiny, text=black!80, inner sep=1pt, align=center},
  ttl/.style={font=\bfseries\scriptsize, anchor=west},
  cv/.style={draw=sheetBlue, fill=sheetBlue!8, font=\ttfamily\tiny, minimum height=4.4mm,
             minimum width=8mm, inner sep=1pt},
  ctx/.style={box, draw=sheetOrange, fill=sheetOrange!10, font=\ttfamily\tiny, inner sep=2pt},
  heap/.style={box, draw=sheetBrown, fill=sheetBrown!12, font=\ttfamily\tiny, inner sep=2pt},
  ptr/.style={->, thick, draw=sheetGrey},
  wk/.style={->, thick, dashed, draw=sheetGreen!70!black},
}

\begin{document}

\sheettitle{Closures in depth — capture, escaping, autoclosure, Sendable}{swift · folio}

\oneliner{A closure is a \textbf{function + a context} holding what it captured, and it is a
\textbf{reference type}. By default it captures \textbf{variables}, not values (a captured
\texttt{var} moves into a shared heap box); a \textbf{capture list} copies values \emph{when the
closure is created}. Parameters are \textbf{non-escaping} by default — \texttt{@escaping} says
the closure may outlive the call, and that is what brings heap contexts, ARC and cycles.}

\vspace{2pt}
\noindent\begin{tikzpicture}[sheet]
  % ── A: shared box ──
  \node[ttl] at (0,3.3) {A \ \texttt{let d = c} — one context, shared};
  \node[lbl, anchor=east] at (0.35,2.6) {\texttt{c}}; \node[cv, anchor=west] (c1) at (0.4,2.6) {fn};
  \node[cv, anchor=west] (c2) at (c1.east) {ctx};
  \node[lbl, anchor=east] at (0.35,1.7) {\texttt{d}}; \node[cv, anchor=west] (d1) at (0.4,1.7) {fn};
  \node[cv, anchor=west] (d2) at (d1.east) {ctx};
  \node[ctx] (cx) at (3.05,2.15) {context};
  \node[heap] (bx) at (4.6,2.15) {box: n = 3};
  \draw[ptr] (c2.east) -- (cx); \draw[ptr] (d2.east) -- (cx); \draw[ptr] (cx) -- (bx);
  \node[lbl, text=sheetBrown] at (2.6,1.05) {\texttt{makeCounter}'s frame is gone —\\
        the escaping capture moved \texttt{n} to the heap};
  \node[lbl] at (2.6,0.35) {\texttt{c(); c(); d()} $\to$ \textbf{3}: copying a closure\\copies the
        \emph{reference}, never the state};
  % ── B: variable vs capture list ──
  \begin{scope}[xshift=57mm]
  \node[ttl] at (0,3.3) {B \ What the context holds};
  \node[heap] (x) at (4.3,2.65) {var x = 99};
  \node[cv, minimum width=17mm] (live) at (0.9,2.65) {\{ print(x) \}};
  \node[ctx] (lc) at (2.45,2.65) {ref to x};
  \draw[ptr] (live) -- (lc); \draw[ptr] (lc) -- (x);
  \node[lbl, anchor=west, text=sheetBrown] at (3.55,2.2) {sees later writes $\to$ 99};
  \node[cv, minimum width=17mm] (snap) at (0.9,1.6) {\{ [x] in \ldots \}};
  \node[ctx] (sc) at (2.45,1.6) {x = 1 (copy)};
  \draw[ptr] (snap) -- (sc);
  \node[lbl, anchor=west] at (3.35,1.6) {copied at \emph{creation} $\to$ 1};
  \node[cv, minimum width=17mm] (wk) at (0.9,0.55) {\{ [weak self] \}};
  \node[ctx] (wc) at (2.45,0.55) {weak self?};
  \node[box, font=\ttfamily\tiny, inner sep=2pt] (vc) at (4.3,0.55) {VC};
  \draw[ptr] (wk) -- (wc); \draw[wk] (wc) -- (vc);
  \node[lbl, anchor=west] at (3.3,0.05) {non-owning; \texttt{nil} once VC dies};
  \end{scope}
  % ── C: lifetimes ──
  \begin{scope}[xshift=112mm]
  \node[ttl] at (0,3.3) {C \ Escaping = outlives the call};
  \draw[thick, sheetBlue] (0.2,2.55) -- (3.1,2.55);
  \draw[thick, sheetBlue] (0.2,2.45) -- (0.2,2.65); \draw[thick, sheetBlue] (3.1,2.45) -- (3.1,2.65);
  \node[lbl, text=sheetBlue] at (0.2,2.85) {\texttt{f(g)} called}; \node[lbl, text=sheetBlue] at (3.1,2.85) {\texttt{f} returns};
  \node[circle, fill=sheetGreen, inner sep=1.5pt] at (1.3,2.55) {};
  \node[lbl, text=sheetGreen!50!black, align=left, anchor=north west] at (0.2,2.3)
        {\textbf{non-escaping} (default): \texttt{g} runs only inside\\the window — context may
         live on the stack, no\\retain, implicit \texttt{self} is fine};
  \draw[thick, sheetOrange, dashed] (3.1,2.55) -- (4.9,2.55);
  \node[circle, fill=sheetOrange, inner sep=1.5pt] at (4.6,2.55) {};
  \node[lbl, text=sheetOrange] at (4.6,2.85) {later};
  \node[lbl, text=sheetOrange!80!black, align=left, anchor=north west] at (0.2,1.15)
        {\textbf{\texttt{@escaping}}: stored / async / \texttt{Task} — heap\\context,
         captures retained $\to$ cycles possible;\\explicit \texttt{self.} required (class)};
  \end{scope}
\end{tikzpicture}

\begin{multicols}{2}
\raggedright

\section{How it works}
\begin{itemize}
  \item \textbf{Capture}: without a list a closure refers to the \emph{variable}; reads and
        writes are shared with the enclosing scope. \textbf{\texttt{[x]}} is an immutable copy made
        when the closure is \emph{created}; for a class it copies the pointer (same object).
        \texttt{[weak self]}, \texttt{[unowned self]}, \texttt{[y = expr]} are the same mechanism.
  \item \textbf{Non-escaping by default} (Swift 3, SE-0103): may only be called during the call,
        cannot be stored or captured by an escaping closure. An \textbf{optional} closure
        parameter \texttt{(() -> Void)?} is \textbf{implicitly escaping}.
  \item \textbf{Implicit \texttt{self}} in escaping closures: allowed when \texttt{self} is a value
        type or listed as \texttt{[self]} (5.3, SE-0269), and after
        \texttt{guard let self} inside \texttt{[weak self]} (5.8, SE-0365).
  \item \textbf{\texttt{withoutActuallyEscaping(f) \{ g in \ldots \}}} lends a non-escaping
        \texttt{f} to an API typed \texttt{@escaping} (e.g. \texttt{lazy.filter}); at the end of the
        block the runtime checks \texttt{g} was not kept — if it was, \textbf{crash}.
  \item \textbf{\texttt{@autoclosure () -> T}}: the caller writes a plain expression, the callee
        decides \emph{whether and when} to evaluate it (\texttt{assert}, \texttt{precondition},
        \op{??}, \op{\&\&}, \op{||}, \texttt{XCTAssert*}, \texttt{@StateObject(wrappedValue:)}).
        Forward it by calling it: \texttt{inner(msg())}. Combine with \texttt{@escaping} to store it.
  \item \textbf{\texttt{@Sendable}} (5.5, SE-0302): may run concurrently, so every capture must
        be \texttt{Sendable} and a local \texttt{var} can't be captured at all (not even read) —
        copy it with \texttt{[x]}. Swift 6: \texttt{Task \{\}} takes a \texttt{sending} closure, so a
        non-Sendable value may go in if nothing uses it afterwards.
  \item \textbf{Function types}: \texttt{(Int) async throws -> Int}; \texttt{rethrows} HOFs
        (\texttt{map}). \texttt{Type.method} is curried \texttt{(Type) -> (Args) -> R};
        \texttt{obj.method} is a closure that \textbf{retains} \texttt{obj}.
  \item \textbf{Trailing closures}: last closure after the parens; \textbf{multiple} (5.3,
        SE-0279): first unlabeled, the rest keep their labels.
\end{itemize}

\section{Example — the syntax ladder}
\begin{lstlisting}[language=SwiftSheet]
names.sorted(by: { (a: String, b: String) -> Bool in
  return a < b })                     // full form
names.sorted(by: { a, b in a < b })   // implicit return
names.sorted { $0 < $1 }             // shorthand + trailing
names.sorted(by: <)                  // operator = function
users.map(\.name)                    // key path (5.2)
UIView.animate(withDuration: 0.3) { v.alpha = 0 }
  completion: { _ in v.removeFromSuperview() }  // 2 trailing
\end{lstlisting}

\columnbreak

\section{Example — capture}
\begin{lstlisting}[language=SwiftSheet]
func makeCounter() -> () -> Int {
  var n = 0                       // boxed: outlives the frame
  return { n += 1; return n } }
let c = makeCounter(); let d = c  // one shared context
c(); c(); print(d())              // 3
var x = 1
let live = { print(x) }           // captures the variable
let snap = { [x] in print(x) }    // copies the value now
x = 99; live(); snap()            // 99, then 1
var fs: [() -> Int] = []
for i in 0..<3 { fs.append { i } }       // new i each pass: 0 1 2
var j = 0
while j < 3 { fs.append { j }; j += 1 }  // one j: 3 3 3
func debugLog(_ m: @autoclosure () -> String) {
  if isDebug { print(m()) } }            // m may never run
\end{lstlisting}

\section{Traps}
\begin{itemize}
  \trap{\textbf{Loop trap}: \texttt{for-in} makes a \emph{fresh} binding per iteration
        ($\to$ 0 1 2); one \texttt{var} mutated by a \texttt{while} is shared ($\to$ 3 3 3).}
  \trap{\textbf{Closures have no identity}: not \texttt{Equatable}, no \op{===} — to remove a
        handler, register it under a token/ID.}
  \trap{\texttt{button.onTap = handleTap} passes \texttt{self.handleTap}: \texttt{self} is
        retained strongly $\to$ cycle. Wrap it: \texttt{\{ [weak self] in self?.handleTap() \}}.}
  \trap{An escaping closure in a struct's \texttt{mutating} method cannot capture
        \texttt{self} (it is \texttt{inout}) — copy the values you need first.}
  \trap{A trailing closure inside an \texttt{if}/\texttt{guard} condition confuses the parser:
        \texttt{if xs.contains(where: \{ \$0 > 3 \}) \{}.}
  \trap{\texttt{@autoclosure} hides laziness from the caller: a side effect in the argument may
        run later, twice, or never.}
  \trap{\texttt{defer} is not a closure: it reads the variable's value at scope exit.}
\end{itemize}

\section{Remember}
\textbf{``Variables by reference, lists by value, escaping to the heap.''}


\end{multicols}

\noindent{\footnotesize\color{sheetGrey}\textit{Related:} ARC \& ownership (retain cycles,
\texttt{weak}/\texttt{unowned}) · Swift 6 strict concurrency (Sendable, \texttt{sending}) ·
operators (\op{\&\&} short-circuit) · key paths as functions · result builders}

\end{document}
